\section{Introduction}\label{sec:v8-introduction}

Researchers increasingly use large language models to measure long
political texts. Once a document exceeds the practical context
limit, the analyst no longer observes the original object: the
input to the final coding step is a compressed surrogate produced by
chunking, retrieval, or recursive summarization. The problem that
follows is one of measurement design, not just prediction. If the
compressed representation drops evidence that depends on combinations
across the document (for example, a tax cut promised in one section
paired with a deficit rule four sections later, where the policy
position turns on the link between the two), the surrogate is acting
on the wrong object, and the downstream label inherits that
distortion no matter how accurate the labeler is on its own.

Design-based supervised learning has formalized the analogous problem
one step downstream. \citet{EgamiEtAl2024} show that even highly
accurate surrogate labels can bias downstream estimates when they
substitute for gold-standard labels without appropriate sampling and
overlap conditions. Their fix is to log a sampling design over the trusted labels and
use it to make the downstream inference valid. Long-document measurement
adds an earlier failure mode: the surrogate is also a compressed
representation of the original document, not just a predicted label.
A retrieval or long-context score may be accurate on average and
still tell the analyst nothing about which cross-document
interactions survived compression.

C-TreePO addresses the representation-level problem. A compression
tree partitions a manifesto into contiguous leaves, summarizes each
leaf into a policy-evidence state, merges sibling states up the tree,
and scores the root. The state map \(g\) is valid for an oracle
\(f^\ast\) when three local laws hold on the realized tree.
Sufficiency: leaf summaries preserve the evidence relevant to the
task. Idempotence: a state stays invariant when passed through the
summarizer again. Merge consistency: a merged state preserves the
parent span it represents. The laws are local in scope and global in
consequence: a valid root state can replace the original manifesto
for any oracle-compatible readout.

The C-TreePO objective has the same two-level shape. The root channel
fits the observed document-level expert target. The local-law channel
trains and audits the state map inside the realized tree. Direct
\(f^\ast\) calls are expensive, so we use a fitted proxy
\(\widehat f\) for the dense local supervision and sample
\(f^\ast\) at a small set of nodes to correct the proxy/oracle gap.

Audit happens at tree nodes, so the analyst picks the unit: a
leaf, a paragraph, a section, or the full platform. Design-based
machinery aggregates the sampled node labels into a document-level
guarantee, and the same design trades one expensive full-document
expert read for several cheaper reads of shorter units at
comparable accuracy. The same supervision design also trains the
compression operators under a fixed audit budget; that is the
optimization side of the framework, the ``PO'' in C-TreePO.

The certificate makes the resulting measurement error reportable
rather than implicit. Section~\ref{sec:v8-audit-certificate}
decomposes the bound on downstream estimate drift into four
interpretable terms. A sampled local-distortion term tells us how
far the produced policy states moved from the raw spans they
represent on the realized tree. A calibration term tracks the gap
between the accessible local judge and the target oracle: the part
of error that more sampled labels cannot remove. A sampling term
carries Horvitz--Thompson uncertainty from the audit design. A
clipping term prices any variance control applied to
inverse-propensity weights. Each term points at a specific lever:
train the state map harder, change the local judge, sample more
nodes, or report the clipping bias explicitly.
Section~\ref{sec:v8-discussion} reads these four terms in the
manifesto application.

We test the framework on the Manifesto/Benoit benchmark
\citep{BenoitEtAl2025}. The benchmark pairs \(235\) held-out party
platforms with expert-survey means on six seven-point policy
dimensions, from economic policy to decentralization. We score each
C-Tree root against the matching expert mean. Per-dimension trees
reach macro \(r=0.829\); the universal-summary tree reaches
\(0.807\). Both clear the matched open-weight baseline at
\(0.793\), and the per-dimension tree clears the proprietary
18-score ensemble at \(0.817\).

\paragraph{Map of the paper and appendices.}
Section~\ref{sec:v8-manifesto-mechanism} introduces the manifesto
policy-state mechanism. Section~\ref{sec:v8-ctree-math} defines
C-Trees and the local laws. Section~\ref{sec:v8-objective} gives the
objective and theorem ladder. Section~\ref{sec:v8-audit-certificate}
gives the audit estimator and finite-sample certificate.
Section~\ref{sec:v8-manifesto-results} reports the application.
Section~\ref{sec:v8-scope-related} places the work against
measurement, long-context, compression, sketching,
preference-learning, and neural-operator literatures.
Section~\ref{sec:v8-discussion} reads the certificate, spends root
and node labels, and names manifesto failure modes.
Appendix~\ref{app:v8-objective-ipw} derives the proxy-corrected
local-law objective.
Appendix~\ref{app:v8-mergeable-nesting} relates classical mergeable
summaries to C-Tree local-law instances.
Appendix~\ref{app:v8-sketch-comparisons} reports sketch comparisons
and classical controls.
Appendix~\ref{app:v8-neural-operator-route} gives the
neural-operator capacity route for the state map \(g\).
Appendix~\ref{app:v8-claim-tiers-manifesto} records claim tiers, the
Benoit replication, and Manifesto provenance.
Appendix~\ref{app:v8-theorem-certificate-details} contains the full
theorem and certificate proofs.
Appendix~\ref{app:v8-manifesto-audit-design} specifies the manifesto
node-audit design.
Appendix~\ref{app:v8-lean-crosswalk} maps the main theorems to their
machine-checked Lean counterparts.
